\begin{abstract}
Augmentative and alternative communication (AAC) devices that leverage large language models can generate context-aware responses tailored to a user's environment, but the computational cost of processing contextual prompts at inference time introduces latency that is incompatible with real-time communication. We present Context-Aware Predictive KV-Cache Priming (CAP-KVC), a system that amortizes this cost by using environmental sensor signals---GPS, Bluetooth Low Energy proximity, and time of day---to predict the user's likely context and pre-compute the corresponding key-value cache tensors during idle periods. At inference time, the pre-computed cache is restored from persistent storage in 2.63~ms on average, and only the user's intent tokens are processed by the language model. Evaluated on a OnePlus~11R with Qwen2.5-0.5B-Instruct (Q4\_K\_M), CAP-KVC reduces end-to-end inference latency by 1.32$\times$ at comparable context length and by 4.59$\times$ as context grows to approximately 500~tokens, relative to standard retrieval-augmented generation. The system is implemented as a complete Android application with a C++/JNI inference backend, a bounded KV cache residency manager supporting three simultaneous states within a 121.6~MB native heap budget, and a multimodal intent classification pipeline combining surface electromyography and gaze tracking through late fusion. We report all measurements, acknowledge limitations including single-device evaluation and synthetic fusion training data, and release the complete implementation for reproducibility.
\end{abstract}
